\documentclass[11pt,a4paper]{article}

% ============================================================
%  ZÁKLADNÍ BALÍČKY
% ============================================================
\usepackage[utf8]{inputenc}
\usepackage[T1]{fontenc}
\usepackage[czech]{babel}       % čeština - dělení slov, uvozovky atd.
\usepackage{geometry}
\geometry{margin=2.5cm}

\usepackage{amsmath}            % rozšířené matematické prostředí
\usepackage{amssymb}            % matematické symboly (\mathbb, \mathcal ...)
\usepackage{amsthm}             % věty, definice, důkazy

\usepackage{xcolor}             % barvy
\usepackage{hyperref}           % klikací odkazy a obsah
\usepackage{enumitem}           % úprava seznamů

\usepackage{bm}                 % tučné matematické symboly
\usepackage{tikz}               % kreslení grafů a stromů

% Pozn.: pokud budeš později chtít obrázky/diagramy stromů,
% přidej \usepackage{tikz} (kreslení) nebo \usepackage{graphicx}
% (vkládání hotových obrázků). Kód, listing apod. zatím nepotřebuješ,
% máš vlastní pseudocode prostředí níž.

% ============================================================
%  VLASTNÍ MATEMATICKÁ MAKRA
% ============================================================

% Množiny čísel
\newcommand{\N}{\mathbb{N}}
\newcommand{\Z}{\mathbb{Z}}
\newcommand{\R}{\mathbb{R}}

\newcommand{\tab}{\hspace*{2em}}

% Časová/prostorová složitost
\newcommand{\bigO}[1]{\mathcal{O}(#1)}
\newcommand{\bigOmega}[1]{\Omega(#1)}
\newcommand{\bigTheta}[1]{\Theta(#1)}

% Zkrácené závorky s automatickou velikostí
\newcommand{\set}[1]{\left\{#1\right\}}
\newcommand{\ceil}[1]{\left\lceil#1\right\rceil}
\newcommand{\floor}[1]{\left\lfloor#1\right\rfloor}

% Prostředí pro věty a definice
\newtheorem{theorem}{Věta}[section]
\newtheorem{lemma}[theorem]{Lemma}
\newtheorem{corollary}[theorem]{Důsledek}
\theoremstyle{definition}
\newtheorem{definition}[theorem]{Definice}
\theoremstyle{remark}
\newtheorem{remark}[theorem]{Poznámka}


% pseudokód
\newcommand{\funkce}[2]{\noindent{\ttfamily\textcolor{blue}{#1}(#2)}}

\newenvironment{pseudocode}
  {\ttfamily\begin{enumerate}[label=(\arabic*), leftmargin=2.2em, itemsep=-0.15em, topsep=0.1em]}
  {\end{enumerate}}

% ============================================================
%  DOKUMENT
% ============================================================
\title{AG1 skripta}
\author{Cyril Ovečka}

\begin{document}

\maketitle
\tableofcontents
\newpage


% ============================================================
% Stromy
% ============================================================

\section{Stromy}
Stromy jsou jedny z nejdůležitějších skupin grafů. Budeme je využívat v algoritmech stejně tak jako v mnoha datových strukturách. V této části si formálně zadefinujeme strom a podíváme se na jeho vlastnosti.

\begin{definition}
  Graf $G$ je \textbf{strom}, právě když je souvislý a acyklický.
\end{definition}

Strom je tedy libovolný souvislý graf, který neobsahuje kružnice.

Pro nesouvislý acyklický graf používáme název \textbf{les}.
\begin{definition}
  Graf $G$ je \textbf{les}, právě když je acyklický (nemusí být souvislý).
\end{definition}

\begin{definition}
  Vrchol $v$ v grafu $G$ nazveme \textbf{list}, pokud platí pokud $\deg_G(v) = 1$.
\end{definition}

List je tedy jakýkoli vrchol v libovolném grafu stupně 1. Ještě si dovolíme poznámku: V anglických textech se list (tedy \textit{leaf}) používá pouze pro vrcholy stupně 1 v lesích nebo stromech. V obecném grafu se nazývají \textit{pendent vertex}.


\vspace{1mm}
\begin{tikzpicture}[
    vnode/.style={circle, draw, minimum size=7mm, inner sep=0pt}
]
    % --- Graph (a) ---
    \node[vnode] (a2) at (0, 2) {2};
    \node[vnode] (a1) at (-1, 0.8) {1};
    \node[vnode] (a4) at (1.1, 2.7) {4};
    \node[vnode] (a5) at (1.5, 1.2) {5};
    \node[vnode] (a3) at (0.7, 0.1) {3};

    \draw (a2) -- (a1);
    \draw (a2) -- (a4);
    \draw (a2) -- (a5);
    \draw (a5) -- (a3);

    \node at (0.2, -0.8) {(a) Strom};

    % --- Graph (b) ---
    \node[vnode] (b2) at (5, 1.8) {2};
    \node[vnode] (b1) at (4.1, 0.8) {1};
    \node[vnode] (b3) at (5.4, 0.1) {3};
    \node[vnode] (b4) at (6.2, 2.7) {4};
    \node[vnode] (b5) at (6.6, 1.2) {5};

    \draw (b2) -- (b1);
    \draw (b2) -- (b3);
    \draw (b4) -- (b5);

    \node at (5.3, -0.8) {(b) Les};
    
    % Bounding box for padding
    \useasboundingbox (-1.5, -1.2) rectangle (8, 3.5);
\end{tikzpicture}
\vspace{0.5cm}

\subsection{Vlastnosti stromu}

\begin{theorem}[Věta o trhání listů]
Každý strom $T$ s alespoň $2$ vrcholy obsahuje alespoň dva listy.
\end{theorem}

\begin{proof}
Důkaz provedeme sporem.

\medskip
\noindent\textbf{Myšlenka důkazu:} Vybereme si nějakou nejdelší cestu ve stromě a tvrdíme, že její krajní vrcholy jsou listy. Pokud si představíte, jak taková cesta vypadá, dává tato myšlenka smysl. Spor získáme jednoduše tak, že tuto nejdelší cestu prodloužíme. Prodlužování nejdelších cest je velmi časté v důkazech sporem a jistě se vám bude ještě hodit.

\medskip
\noindent\textbf{Důkaz:} Pro spor nechť je graf $G$ strom a $P$ je nějaká nejdelší cesta v grafu $G$ s krajními vrcholy $u$ a $v$. Předpokládejme, že alespoň jeden z nich není list v grafu $G$.

BUNO $u$ je krajní vrchol, který není list. Musí tedy mít alespoň $2$ hrany (nemůže mít $0$, protože je součástí cesty). Označme $e = \{u,w\}$ hranu, která není obsažena v $P$. Nyní rozlišujeme dva případy:

\begin{enumerate}
    \item Pokud by $w$ byl součástí cesty $P$, pak by v grafu byl cyklus. To je spor s tím, že $G$ je strom (strom je acyklický).
    \item Pokud by $w$ součástí cesty nebyl, pak lze cestu $P$ prodloužit přidáním vrcholu $w$ za vrchol $u$. To je spor s předpokladem, že $P$ je nejdelší cesta v $G$.
\end{enumerate}

V obou případech jsme dospěli ke sporu. Tvrzení tedy platí.
\end{proof}

 


% ============================================================
% Binární vyhledávací strom
% ============================================================

\section{Binární vyhledávací strom}

\begin{definition}
Strom nazveme \textbf{binární}, pokud
    \begin{itemize}
        \item je zakořeněný
        \item každý vrchol má nejvýše dva syny
        \item u synů rozlišujeme, který je levý a který pravý
    \end{itemize}
\end{definition} \bigskip

\begin{remark}
Zakořeněnění je důležitá vlastnost. Představte si, že dostanete binární vyhledávací strom a prvek co se má do něj vložit.
Kde začnete? Bez zakořenění neexistuje vstupní bod.
\end{remark} \bigskip

\begin{definition}[Značení]
Pro vrchol $v$ binárního stromu $T$ značíme
\begin{itemize}
    \item $\ell(v)$ a $r(v)$: levý a pravý syn vrcholu $v$,
    \item $L(v)$ a $R(v)$: levý a pravý podstrom vrcholu $v$,
    \item $p(v)$: otec vrcholu $v$,
    \item $T(v)$: podstrom tvořený kořenem $v$ a všemi jeho potomky,
    \item $h(T(v))$: hloubka stromu $T(v)$ je počet hladin $T(v)$,
    \item $|T|$, $|T(v)|$, $|L(v)|$ a $|R(v)|$: počet vrcholů stromu $T$, $T(v)$, $L(v)$ a $R(v)$.
\end{itemize}
Pokud vrchol $v$ nemá levého syna, položíme $\ell(v) = \emptyset$. Podobně pro $r(v)$ a $p(v)$.\newline
Pak dodefinujeme, že $T(\emptyset)$ je prázdný strom a $h(T(\emptyset)) = 0$.
\end{definition} \bigskip

\noindent Nyní si zavedeme binární vyhledávací strom. Kdykoliv se budeme bavit o binárních stromech je myšlen tento.

\begin{definition}
    Binární vyhledávací strom (BVS) je binární strom, v jehož každém vrcholu v je uložen \textbf{unikátní} klíč k(v) a pro každý jeho vrchol platí:

    \begin{itemize}
        \item Pokud $a \in L(v)$, pak $k(a) \bm{<} k(v)$
        \item Pokud $b \in R(v)$, pak $k(b) \bm{<} k(v)$
    \end{itemize}
\end{definition} \bigskip

\begin{remark}
Na BVS je pár důležitých operací. Doporučuju pochopit princip těchto algoritmů než se učit pseudokód nazpaměť.
\end{remark} \bigskip

\noindent První je algoritmus BVSShow. Jeho cílem je provést InOrder průchod stromu a vypsat všechny klíče vzestupně. Z definice BVS víme, že pro každý vrchol $v$ platí: klíče v jeho levém podstromě jsou menší a klíče v pravém podstromě větší. Aby algoritmus vypsal klíče ve vzestupném pořadí, musí se nejprve co nejvíce zanořit doleva, k nejmenším klíčům. Teprve po návratu z levého syna vypíše klíč aktuálního vrcholu a poté pokračuje do pravého syna, kde jsou klíče větší. Tento postup se opakuje pro každý vrchol stromu.
Jelikož v každém vrcholu strávíme konstantní čas, celková časová složitost algoritmu je $\bigO{|T(v)|}$. \bigskip

\noindent \textbf{Vstup}: Ukazatel na kořen $v$ libovolného BVS $T(v)$ \newline
\noindent \textbf{Výstup}: Vzestupně uspořádaná posloupnost klíčů všech vrcholů v $T(v)$ \newline
\funkce{BVSShow}{ukazatel na vrchol $v$}
\begin{pseudocode}
  \item Pokud $v = \emptyset$: return \textcolor{gray}{//$T(v)$ je prázdný}
  \item Zavoláme \textcolor{blue}{BVSShow}($\ell(v)$)
  \item Vypíšeme $k(v)$
  \item Zavoláme \textcolor{blue}{BVSShow}($r(v)$)
\end{pseudocode} \bigskip

\noindent BVSMin/BVSMax. Jeho cílem je vypsat minimální/maximální hodnotu klíče ve stromě. Pro minimální klíč se z daného vrcholu zanořujeme stále doleva. Jakmile vrchol již levého syna nemá, pak je minimálním. Analogicky pro maximum, pouze se zanořujeme doprava. V kažédm vrcholu provedeme konstantě operací tedy celkem provedeme $\bigO{h(T(v))}$\bigskip

\noindent \textbf{Vstup}: Ukazatel na vrchol $v$ libovolného BVS $T(v)$ \newline
\noindent \textbf{Výstup}: Minimimální klíč daného BVS $T(v)$ \newline
\funkce{BVSMin}{ukazatel na vrchol $v$}
\begin{pseudocode}
  \item Pokud $v = \emptyset$: return \textcolor{gray}{//$T(v)$ je prázdný}
  \item Pokud $\ell(v) = \emptyset$: return $k(v)$
  \item Zavoláme \textcolor{blue}{BVSMin}($\ell(v)$)
\end{pseudocode} \bigskip

\noindent BVSFind slouží k nalezení vrcholu ke klíči v BVS. Dostaneme na vstupu ukazatel na kořen $v$ BVS a hledaný klíč. Nyní postupuje následovně, pokud je hledaný klíč větší než $k(v)$. Tak musíme jít do pravého podstromu $R(v)$, kde jsou klíče větší. Analogicky pro menší klíč. Dokud buďto nenalezneme vrchol s hledaným klíčem nebo nevylezeme ze stromu. V každém vrcholu provedeme pouze konstantě mnoho operací a celkem tedy $\bigO{h(T(v))}$. \bigskip

\noindent \textbf{Vstup}: Ukazatel na vrchol $v$ libovolného BVS $T(v)$ a hledaný klíč $x$. \newline
\noindent \textbf{Výstup}: Ukazatel na vrchol s klíčem $x$, pokud takový v $T(v)$
existuje, jinak $\emptyset$. \newline
\funkce{BVSFind}{ukazatel na vrchol $v$ a hledaný klíč $x$}
\begin{pseudocode}
    \item Pokud $v = \emptyset$: return $\emptyset$
    \item Pokud $k(v) > x$: return \textcolor{blue}{BVSFind}($\ell(v)$)
    \item Pokud $k(v) < x$: return \textcolor{blue}{BVSFind}($r(v)$)
    \item return $v$
\end{pseudocode} \bigskip

\noindent BVSPred a BVSSucc slouží k získání předchůdce/následníka nějakého vrcholu $v$ v BVS $T$. Pro předchůdce, je to vždy takový klíč $x$, který je největší ze všech klíčů menších než $x$ v $T$. Jak ho nalezneme? Rozlišujeme zde dva případy, ten lehčí je když vrchol $v$ má levého syna. V tom případě je předchůdce vrchol s maximálním klíčem z celého levého podstromu $L(v)$. To platí, jelikož jsou všechny klíče v $L(v)$ menší než $k(v)$ a maximum z $L(v)$ je ten největší menší. Druhý případ je když $v$ nemá levého syna v tom případě se musíme vynořit nahoru (doprava jít nemůžeme tam jsou všechny klíče větší). Takto cestujeme směrem ke kořeni až do doby než se vrátíme do rodiče $p$ zprava. Předchůdce je poté vrchol $p$. Zprava se chceme vrátit jelikož my potřebujeme najít menší první menší klíč. Pokud by jsme přišli zleva, tak je rodič větší než celý náš podstrom.\newline
Pokud náš prvek je minimálním v celém stromě vrátíme nic. Alternativně lze dodefinovat vrcholy s klíčy $\pm\infty$.\newline
Analogicky platí i pro následovníka. Stačí prohodit slova levý a pravý.
buďto využijeme BVSMax nebo procházíme po rodičích a případně do kořene. celkem tedy za $\bigO{h(T(v))}$. \bigskip

\noindent \textbf{Vstup}: Ukazatel na kořen $v$ libovolného BVS $T(v)$ a nějaký jeho vrchol $w$. \newline
\noindent \textbf{Výstup}: Ukazatel na předchůdce vrcholu $w$ v BVS $T$.
existuje-li, jinak $\emptyset$. \newline
\funkce{BVSPred}{ukazatel na kořen $v$ a vrchol $w$}
\begin{pseudocode}
    \item Pokud $l(w) \neq \emptyset$
    \item \tab return \textcolor{blue}{BVSMax}($l(w)$) \textcolor{gray}{//Máme levého syna, tedy přechůdce je maximem $L(w)$}
    \item $z := p(w)$
    \item Dokud ($z \neq \emptyset$ $\And$ $w = l(z)$) \textcolor{gray}{//Jdeme směrem ke kořeni a čekáme až přijdeme zprava}
    \item \tab $w := z$
    \item \tab $z := p(z)$
    \item return $z$
\end{pseudocode} \bigskip

\noindent BVSInsert jak z názvu plyne vkládá prvek do BVS. Na vstupu dostaneme ukazatel na kořen $v$ nějakého BVS a klíč $x$, který chceme vložit. Ve stromě předpokládáme unikátní klíče v případě již existujícího stejného klíče se vložení neprovede. Algoritmus prochází strom stejně jako BVSFind a hledá volné místo kam vložit prvek. Jakmile se dostaneme mimo strom, tak jsme nalezli místo kam vložit daný klíč. Pokud je strom prázdný, tak pouze vložíme prvek. V každém vrcholu provedeme pouze konstantě mnoho operací a jdeme skrze vrcholy od kořene až do listu, což je celkem za $\bigO{h(T(v))}$. \bigskip

\noindent \textbf{Vstup}: Ukazatel na kořen $v$ libovolného BVS $T(v)$ a vkládaný klíč $x$. \newline
\noindent \textbf{Výstup}: Ukazatel na kořen BVS s klíčem $x$. \newline
\funkce{BVSInsert}{Ukazatel na kořen $v$ BVS a vkládaný klíč $x$}
\begin{pseudocode}
    \item Pokud $v = \emptyset$: \textcolor{gray}{//Strom je prázdný}
    \item \tab Vytvoř nový vrchol $v$
    \item \tab return $v$in overleafin overleaf
    \item Pokud $k(v) > x$:
    \item \tab $l(v) :=$ \textcolor{blue}{BVSInsert}$(l(v), x)$
    \item Pokud $k(v) < x$:
    \item \tab $r(v) :=$ \textcolor{blue}{BVSInsert}$(r(v), x)$
    \item return $v$
\end{pseudocode} \bigskip

\noindent Poslední operací je BVSDelete, ten slouží k odstranění vrcholu $v$ z BVS dle zadaného klíče. Zde máme 3 možnosti. První je pokud vrchol $v$ nemá žádného syna neboli je listem. V tom případě ho prostě odstraníme a vše funguje dále. Další případ je kdy má právě jednoho syna vrchol $w$. Ani zde neni problém, rodiči vrcholu $v$, $p(v)$ nahradíme syna $v$ za vrchol $w$. Vrchol $v$ poté bezproblému odstraníme. Posledním případem je kdy vrchol $v$ má oba syny. V tomto případě nemůžeme vrchol $v$ odstranit jinak by jsme rozbili strukturu. Pomůžeme si trikem a to pomocí nahrazení vrcholu $v$ za jeho následníka nebo předchůdce (ti musí existovat jelikož má vrchol $v$ oba syny a zároveň jsou to maxima/minima těch podstromů, takže problém se nerekurzí dále). Také jsme nemohli rozbít definici BVS, jelikož jsme nahradili bezprostředním následníkem/předchůdcem a stále jsou vlevo menší klíče a vpravo větší. Hledáme vrchol jako při BVSFind, k tomu hledáme Minimum podstromu pomocí BVSMin a v každém vrcholu provedeme konstantě mnoho operací tedy celkem $\bigO{h(T(v)}$. \bigskip

\noindent \textbf{Vstup}: Ukazatel na kořen libovolného BVS $T(v)$ a klíč $x$. \newline
\noindent \textbf{Výstup}: kazatel na kořen BVS, ze kterého byl odstraněn vrchol s
klíčem $x$, existoval-li. \newline
\funkce{BVSDelete}{Ukazatel na kořen $v$ BVS a mazaný klíč $x$}
\begin{pseudocode}
    \item Pokud $v = \emptyset$: return $\emptyset$
    \item Pokud $k(v) > x$: 
    \item \tab $l(v) := $ \textcolor{blue}{BVSDelete}($\ell(v)$)
    \item Pokud $k(v) < x$:
    \item \tab $r(v) := $ \textcolor{blue}{BVSDelete}($r(v)$)
    \item Pokud $k(v) = x$:
    \item \tab Pokud $\ell(v) = r(v) = \emptyset$: return $\emptyset$ \textcolor{gray}{//Vrchol $v$ je listem, mazání je triviální}
    \item \tab Pokud $\ell(v) = \emptyset$: return $r(v)$ \textcolor{gray}{//Vrchol v má pouze jednoho syna, mazání je triviální}
    \item \tab Pokud $r(v) = \emptyset$: return $\ell(v)$
    \item \tab $w :=$ \textcolor{blue}{BVSMin}($r(v)$)
    \item \tab $k(V) := k(w)$
    \item \tab $r(v) := $ \textcolor{blue}{BVSDelete}($r(v), k(w)$) \textcolor{gray}{//Musíme smazat ještě ten prohozený vrchol}
    \item return $v$
\end{pseudocode} \bigskip

\begin{remark}
BVS nemá jednoznačnou strukturu. Skutečne zkuste provést BVSInsert na prvcích \{2, 1, 3\} a poté na \{1, 2, 3\}. Množiny jsou v obou případech stejné. Ale výsledný BVS je jiný.\bigskip
\end{remark}

\noindent Časová složitost a problém BVS. Již jsme určili že operace jako BVSInsert, BVSDelete, BVSMin, BVSFind, atd. mají časovou složitost $\bigO{h(T(v)}$. Zde ale nastává problém, pokud nám strom zdegraduje na spojový seznam. To se stane např. pokud všechny prvky budou mít pouze pravého syna. Poté je $\bigO{h(T(v)} = \bigO{|T(v)|} = \bigO{n}$. Což se nám nelíbí, jelikož ztrácíme výhody oproti běznému poli. Jistě lze namítnout, že šance na tento scénář je malá. V případě kdy přijde útočník a začne nám dávat taková data, která toto způsobí. tak jsme nahraní. Na ochranu proti takovémuto zdegenerování se zavádí pojmy vyváženosti. \bigskip

\begin{definition}
BVS nazveme \textbf{dokonale vyvážený}, pokud pro každý jeho vrchol $v$ platí:
\[ \left| |L(v)| - |R(v)| \right| \leq 1 \]
\end{definition} \bigskip

\begin{remark}
Dokonale vyvážený BVS o velikosti $n$ má $1 + \floor{\log n}$ hladin. Operace na takovém stromě mají tedy optimální časovou složitost $\bigO{\log n}$.
\end{remark} \bigskip

\noindent Dokonalé vyvážení nám krásně zajišťuje logaritmicou hloubku BVS. Bohužel se ukazuje, že udržet tuto strukturu stromu současně s přidáváním a mazáním prvků, je příliš nákladné. Alespoň jedna z operací BVSInsert a BVSDelete zdegradují na lineární složitost $\bigO{n}$. \bigskip

\begin{theorem}[O složitosti operací nad dokonale vyváženým BVS]
Pokud má BVS zůstávat po provedení operací BVSInsert a BVSDelete dokonale vyvážený, potom ať použijeme jakoukoli implementaci zmíněných operací, bude časová složitost aspoň jedné z nich $\bigOmega{n}$ pro nekonečně mnoho různých $n$.
\end{theorem} \bigskip

\noindent Myšlenka důkazu. Stačí nám ukázat že existuje dokonale vyvážený BVS a nějaká posloupnost operací BVSInsert a BVSDelete kde alespoň jedna bude mít složitost $\bigO{n}$. \bigskip

\begin{proof}
    
\end{proof}

% ============================================================
% AVL Strom
% ============================================================

\section{Avl strom}


\end{document}